\documentclass[preprint,12pt]{elsarticle}

\usepackage{amsmath,amssymb}
\usepackage{graphicx}
\usepackage{booktabs}
\usepackage{longtable}
\usepackage{multirow}
\usepackage{hyperref}
\usepackage{xcolor}
\usepackage[margin=2.5cm]{geometry}

\journal{Computers, Environment and Urban Systems}

\begin{document}

\begin{frontmatter}

\title{Supplementary Material for:\\
Spatiotemporal Analysis of Urban Traffic Congestion, Network Topology, and Capacity Constraints in Indonesian Metropolitan Cities}

\author[undip]{Firman Hadi\corref{cor1}}
\ead{firmanhadi21@lecturer.undip.ac.id}
\author[undip]{Yasser Wahyuddin}
\author[undip]{L.M. Sabri}
\author[nca]{Agung Indrajit}

\affiliation[undip]{organization={Department of Geodetic Engineering, Universitas Diponegoro},
  city={Semarang}, country={Indonesia}}
\affiliation[nca]{organization={Deputy for Green and Digital Transformation, Nusantara Capital Authority},
  country={Indonesia}}

\cortext[cor1]{Corresponding author.}

\end{frontmatter}

\setcounter{section}{0}
\renewcommand{\thesection}{S\arabic{section}}
\renewcommand{\thetable}{S\arabic{table}}
\renewcommand{\thefigure}{S\arabic{figure}}

%% ============================================================
\section{Spatial Weight Sensitivity Analysis}
\label{sec:s1}

Following the OSM-based re-aggregation, we computed Moran's I under five spatial weight specifications for each city: K-nearest neighbours with $k \in \{4, 8, 12\}$ and distance-band weights at 500\,m and 1000\,m thresholds. All specifications confirm strong, statistically significant positive spatial autocorrelation.

\begin{table}[!htbp]
\centering
\caption{Moran's I under alternative spatial weight specifications (evening peak mean jam factor). All 15 city--weight combinations are highly significant ($p \approx 0$), confirming robust positive spatial autocorrelation across all weight specifications.}
\label{tab:s1}
\begin{tabular}{lccccc}
\toprule
City & KNN $k{=}4$ & KNN $k{=}8$ & KNN $k{=}12$ & Dist 500\,m & Dist 1000\,m \\
\midrule
Semarang & 0.347*** & 0.298*** & 0.270*** & 0.323*** & 0.288*** \\
Bandung  & 0.355*** & 0.291*** & 0.248*** & 0.378*** & 0.309*** \\
Jakarta  & 0.368*** & 0.310*** & 0.270*** & 0.362*** & 0.276*** \\
\bottomrule
\end{tabular}

\smallskip
\noindent\textit{*** $p < 0.001$ for all entries. $n$: Semarang = 924, Bandung = 2,583, Jakarta = 13,201 segments.}
\end{table}

Moran's I values are remarkably consistent across weight specifications (range 0.248--0.378), indicating robust spatial autocorrelation that is not an artefact of a particular neighbourhood definition. The distance-band weights at 500\,m and 1000\,m yield comparable values to KNN specifications, confirming that the spatial structure persists at both local and broader spatial scales.

%% ============================================================
\section{Period-Specific Moran's I}
\label{sec:s2}

Global Moran's I was computed separately for each of the eight temporal periods to assess whether spatial clustering varies by time of day. Spatial weights use KNN with $k{=}8$.

\begin{table}[!htbp]
\centering
\caption{Moran's I by temporal period (KNN $k{=}8$). All 24 period--city combinations are highly significant ($p \approx 0$), confirming ubiquitous positive spatial autocorrelation across all times of day. Moran's I increases from night to evening peak, mirroring the temporal congestion gradient.}
\label{tab:s2}
\begin{tabular}{lccc}
\toprule
Period & Semarang & Bandung & Jakarta \\
\midrule
Night             & 0.176*** & 0.196*** & 0.222*** \\
Morning Peak      & 0.289*** & 0.224*** & 0.257*** \\
Morning Off-Peak  & 0.280*** & 0.276*** & 0.280*** \\
Lunch Hours       & 0.241*** & 0.282*** & 0.294*** \\
Afternoon Off-Peak& 0.266*** & 0.282*** & 0.306*** \\
Evening Peak      & 0.298*** & 0.291*** & 0.310*** \\
Evening Off-Peak  & 0.285*** & 0.265*** & 0.284*** \\
Late Night        & 0.255*** & 0.222*** & 0.260*** \\
\bottomrule
\end{tabular}

\smallskip
\noindent\textit{*** $p < 0.001$ for all entries.}
\end{table}

Moran's I ranges from 0.176 (Semarang, night) to 0.310 (Jakarta, evening peak), with a clear temporal gradient: spatial autocorrelation strengthens during peak hours and weakens at night. This pattern suggests that peak-hour congestion is more spatially structured than off-peak congestion, likely reflecting the concentration of commuting flows along predictable corridors.

%% ============================================================
\section{Getis-Ord Gi* Hot and Cold Spot Counts}
\label{sec:s3}

Getis-Ord Gi* statistics were computed with KNN ($k{=}8$) binary weights and 999 permutations to triangulate the LISA findings reported in the main text.

\begin{table}[!htbp]
\centering
\caption{Getis-Ord Gi* hot and cold spot counts (evening peak, $p < 0.05$). Approximately 36\% of segments are classified as significant hot or cold spots, with cold spots outnumbering hot spots in all three cities.}
\label{tab:s3}
\begin{tabular}{lcccc}
\toprule
City & Hot Spots & Cold Spots & Not Significant & \% Significant \\
\midrule
Semarang & 116 (12.6\%) & 222 (24.0\%) & 586 (63.4\%)   & 36.6\% \\
Bandung  & 392 (15.2\%) & 518 (20.1\%) & 1,673 (64.8\%) & 35.2\% \\
Jakarta  & 2,063 (15.6\%) & 2,644 (20.0\%) & 8,494 (64.3\%) & 35.7\% \\
\bottomrule
\end{tabular}
\end{table}

Among hot spots, 1,037 (50.3\%) in Jakarta, 170 (43.4\%) in Bandung, and 47 (40.5\%) in Semarang reach the stricter $p < 0.01$ threshold, indicating robust clustering. The overall proportion of significant segments ($\sim$35--37\%) is substantially higher than the $\sim$5\% expected by chance, and consistent with the strong Moran's I values reported in Section~\ref{sec:s1}.

%% ============================================================
\section{H3 Hexagonal Aggregation: Extended Results}
\label{sec:s4}

This section provides extended results from the H3 hexagonal scale-sensitivity analysis described in Section 4.4 and reported in Section 5.6 of the main text. The Uber H3 hierarchical geospatial indexing system (Brodsky, 2018) was used to aggregate road-segment-level variables into regular hexagonal bins at two resolutions, testing whether the null correlations observed at the segment level reflect genuine absence of relationships or a modifiable areal unit problem (MAUP) artefact.

\subsection{Hexagon Coverage Statistics}

\begin{table}[!htbp]
\centering
\caption{H3 hexagon coverage by city and resolution. Resolution 8 ($\approx$461\,m diameter) captures neighbourhood-scale patterns; resolution 9 ($\approx$174\,m diameter) approximates block-level aggregation.}
\label{tab:s4}
\begin{tabular}{llrrr}
\toprule
City & Resolution & Hexagons & Mean Segments/Hex & Mean Observations/Hex \\
\midrule
Semarang & 8 ($\approx$461\,m) & 292 & 3.2 & 52,065 \\
Semarang & 9 ($\approx$174\,m) & 572 & 1.6 & 26,573 \\
Bandung  & 8 ($\approx$461\,m) & 571 & 4.5 & 75,969 \\
Bandung  & 9 ($\approx$174\,m) & 1,415 & 1.8 & 30,680 \\
Jakarta  & 8 ($\approx$461\,m) & 2,000 & 6.6 & 103,019 \\
Jakarta  & 9 ($\approx$174\,m) & 6,202 & 2.1 & 33,274 \\
\bottomrule
\end{tabular}
\end{table}

The higher segment density per hexagon in Jakarta reflects its denser road network. At resolution 8, each hexagon contains sufficient observations (52,000--103,000) for stable mean estimates.

\subsection{POI and Centrality Correlations Across Spatial Scales}

\begin{table}[!htbp]
\centering
\caption{Spearman correlations ($\rho$) between hex-level jam factor and spatial predictors at two H3 resolutions. Several correlations reach statistical significance at the hexagonal scale, particularly for Bandung and Semarang, though effect sizes remain modest ($|\rho| < 0.20$).}
\label{tab:s5}
\begin{tabular}{llccccc}
\toprule
 & & \multicolumn{2}{c}{POI Density} & \multicolumn{2}{c}{Betweenness Centrality} \\
\cmidrule(lr){3-4} \cmidrule(lr){5-6}
City & Scale & $\rho$ & $p$-value & $\rho$ & $p$-value \\
\midrule
\multirow{2}{*}{Semarang}
 & H3 Res.\ 9 & $-0.064$ & 0.125 & $+0.098$ & 0.019* \\
 & H3 Res.\ 8 & $-0.193$ & 0.001** & $-0.010$ & 0.865 \\
\midrule
\multirow{2}{*}{Bandung}
 & H3 Res.\ 9 & $+0.102$ & $<$0.001** & $+0.093$ & $<$0.001** \\
 & H3 Res.\ 8 & $+0.131$ & 0.002** & $+0.153$ & $<$0.001** \\
\midrule
\multirow{2}{*}{Jakarta}
 & H3 Res.\ 9 & $+0.032$ & 0.012* & $+0.071$ & $<$0.001** \\
 & H3 Res.\ 8 & $+0.018$ & 0.423 & $+0.072$ & 0.001** \\
\bottomrule
\end{tabular}

\smallskip
\noindent\textit{* $p < 0.05$; ** $p < 0.01$}
\end{table}

At the neighbourhood scale (resolution~8), Semarang shows a significant negative POI correlation ($\rho = -0.193$, $p = 0.001$), while Bandung exhibits the strongest positive associations for both POI density ($\rho = +0.131$) and centrality ($\rho = +0.153$). Jakarta's centrality correlation reaches significance ($\rho = +0.072$, $p = 0.001$) but remains practically negligible. These results suggest that the relationship between spatial predictors and congestion operates at the neighbourhood rather than segment scale, though effect sizes are uniformly modest.

\subsection{Global Moran's I Across Spatial Scales}

\begin{table}[!htbp]
\centering
\caption{Global Moran's I for jam factor at three spatial scales. All three cities exhibit strong, statistically significant positive spatial autocorrelation at both the segment and hexagonal levels, with Moran's I increasing at coarser hexagonal resolutions due to smoothing of within-cell noise.}
\label{tab:s6}
\begin{tabular}{llcc}
\toprule
City & Scale & Moran's I & $p$-value \\
\midrule
\multirow{3}{*}{Semarang}
 & Segment (KNN $k{=}8$)    & $+0.298$ & $\approx 0$*** \\
 & H3 Res.\ 9 & $+0.345$ & 0.001*** \\
 & H3 Res.\ 8 & $+0.486$ & 0.001*** \\
\midrule
\multirow{3}{*}{Bandung}
 & Segment (KNN $k{=}8$)    & $+0.291$ & $\approx 0$*** \\
 & H3 Res.\ 9 & $+0.351$ & 0.001*** \\
 & H3 Res.\ 8 & $+0.227$ & 0.001*** \\
\midrule
\multirow{3}{*}{Jakarta}
 & Segment (KNN $k{=}8$)    & $+0.310$ & $\approx 0$*** \\
 & H3 Res.\ 9 & $+0.314$ & 0.001*** \\
 & H3 Res.\ 8 & $+0.247$ & 0.001*** \\
\bottomrule
\end{tabular}

\smallskip
\noindent\textit{*** $p < 0.001$}
\end{table}

Strong spatial autocorrelation is now ubiquitous across all cities and scales (all $I > 0.22$, all $p = 0.001$). Semarang shows the strongest clustering at resolution~8 ($I = 0.486$), suggesting that neighbourhood-level aggregation reveals coherent congestion zones that are partially masked by segment-level heterogeneity. The consistent Moran's I values across scales confirm that congestion is genuinely spatially structured: segments and hexagons with elevated congestion tend to be surrounded by similarly congested neighbours. Importantly, this clustering is largely \emph{not} explained by POI density or betweenness centrality (Table~\ref{tab:s5}), suggesting that other spatial factors---such as land-use mix, population density, and public transit access---drive neighbourhood-level congestion patterns.

%% ============================================================
\section{Spatial Match Quality}
\label{sec:s5}

The spatial join between HERE Traffic segments and OSMnx road network edges was validated by examining centroid-to-centroid match distances.

\begin{table}[!htbp]
\centering
\caption{Spatial match quality: HERE Traffic segments to OSM reference network. With OSM-based aggregation using composite segment IDs, median match distances are effectively zero, and $>$99\% of segments match within 50\,m.}
\label{tab:s7}
\begin{tabular}{lccccc}
\toprule
City & $n$ & Median (m) & \% within 50\,m & \% within 100\,m & \% within 200\,m \\
\midrule
Semarang & 1,225  & 0.0 & 99.3\% & 99.6\% & 99.8\% \\
Bandung  & 3,404  & 0.0 & 99.5\% & 99.8\% & 100.0\% \\
Jakarta  & 17,398 & 0.0 & 99.7\% & 99.8\% & 99.9\% \\
\bottomrule
\end{tabular}
\end{table}

The near-perfect match quality reflects the OSM-based aggregation pipeline, where HERE Traffic segments are directly matched to OSM reference geometries using composite segment identifiers derived from the OSM network. This eliminates the positional uncertainty inherent in the previous centroid-to-centroid spatial join approach.

%% ============================================================
\section{Spatial Regression Results}
\label{sec:s6}

To properly account for spatial dependence, Spatial Lag and Spatial Error models (via PySAL's \texttt{spreg}) were estimated alongside OLS multiple regression. The dependent variable is mean evening-peak jam factor; independent variables are edge betweenness centrality and POI density within a 500\,m buffer.

\begin{table}[!htbp]
\centering
\caption{Spatial regression results: jam factor $\sim$ betweenness centrality $+$ POI density. Across all model specifications, $R^2$ values are negligible ($<$0.3\%), confirming that spatial predictors explain virtually none of the congestion variance.}
\label{tab:s8}
\small
\begin{tabular}{llcccccl}
\toprule
City & Model & $R^2$ & $\beta$ (betw.) & $p$ & $\beta$ (POI) & $p$ & Spatial param. \\
\midrule
\multirow{3}{*}{Jakarta}
 & OLS           & 0.003\% & $-0.048$ & 0.635 & 0.000010  & 0.704 & ---              \\
 & Spatial Lag   & 0.006\% & $-0.337$ & 0.168 & $-0.000005$ & 0.899 & $\rho = 2.731$   \\
 & Spatial Error & 0.003\% & $-0.051$ & 0.615 & 0.000009  & 0.708 & $\lambda = 0.012$ \\
\midrule
\multirow{3}{*}{Bandung}
 & OLS           & 0.260\% & 0.270    & 0.005 & $-0.000010$ & 0.803 & ---              \\
 & Spatial Lag   & 0.128\% & 0.231    & 0.064 & $-0.000000$ & 0.996 & $\rho = 0.350$   \\
 & Spatial Error & 0.260\% & 0.269    & 0.005 & $-0.000010$ & 0.817 & $\lambda = 0.023$ \\
\midrule
\multirow{3}{*}{Semarang}
 & OLS           & 0.004\% & $-0.017$ & 0.894 & $-0.000007$ & 0.887 & ---              \\
 & Spatial Lag   & 0.013\% & 0.033    & 0.866 & 0.000055  & 0.706 & $\rho = 2.568$   \\
 & Spatial Error & 0.004\% & $-0.018$ & 0.886 & $-0.000008$ & 0.867 & $\lambda = -0.034$ \\
\bottomrule
\end{tabular}
\end{table}

Bandung's OLS betweenness coefficient is statistically significant ($p{=}0.005$) but trivially small in magnitude ($R^2 = 0.26\%$). This significance vanishes in the Spatial Lag model ($p{=}0.064$), suggesting it is partially attributable to spatial dependence rather than a genuine centrality effect. The spatial autoregressive parameters ($\rho$, $\lambda$) are small, consistent with the weak global spatial autocorrelation reported in the main text and Table~\ref{tab:s1}.

%% ============================================================
\section{Supplementary Figures}
\label{sec:sfigs}

\subsection{Hotspot Maps}

Figures~\ref{fig:s1}--\ref{fig:s2} show LISA cluster and Getis-Ord Gi* hotspot maps for the evening peak period for Bandung and Semarang. The main text includes the Jakarta hotspot map.

\begin{figure}[!htbp]
\centering
\includegraphics[width=\textwidth]{figures/bdg_hotspots_evening_peak.png}
\caption{LISA cluster and Getis-Ord Gi* hotspot maps --- Bandung (evening peak). Hot spots are concentrated along the northern industrial corridor (Jl.\ Soekarno-Hatta) and the Pasteur--Cipaganti commercial strip.}
\label{fig:s1}
\end{figure}

\begin{figure}[!htbp]
\centering
\includegraphics[width=\textwidth]{figures/smg_hotspots_evening_peak.png}
\caption{LISA cluster and Getis-Ord Gi* hotspot maps --- Semarang (evening peak). Despite the smallest network ($n{=}1{,}076$), identifiable hot spots cluster around the Simpang Lima central roundabout and Jl.\ Pandanaran commercial corridor.}
\label{fig:s2}
\end{figure}

\subsection{H3 Hexagonal Aggregation Figures}

\begin{figure}[!htbp]
\centering
\includegraphics[width=0.75\textwidth]{figures/h3_morans_scale.png}
\caption{Global Moran's I by spatial scale for three cities. Bars show Moran's I at the road-segment level, H3 resolution~9 ($\approx$174\,m), and H3 resolution~8 ($\approx$461\,m). All values are highly significant ($p = 0.001$). Spatial autocorrelation is strong at all scales (all $I > 0.22$), with Semarang showing the strongest clustering at the neighbourhood scale ($I = 0.486$ at resolution~8).}
\label{fig:s3}
\end{figure}

\begin{figure}[!htbp]
\centering
\includegraphics[width=\textwidth]{figures/h3_correlation_panels.png}
\caption{Scatter plots of POI density (top row) and network centrality (bottom row) versus hex-level mean jam factor at H3 resolution~8. Dashed lines show OLS regression fits. Pearson $r$ and Spearman $\rho$ are annotated per panel. Bandung shows the strongest correlations (POI $\rho = +0.131$, centrality $\rho = +0.153$), though effect sizes remain modest.}
\label{fig:s4}
\end{figure}

\begin{figure}[!htbp]
\centering
\includegraphics[width=0.85\textwidth]{figures/h3_correlation_comparison.png}
\caption{Spearman $\rho$ for POI density (left) and network centrality (right) at three spatial scales: road segment, H3 resolution~9 ($\approx$174\,m), and H3 resolution~8 ($\approx$461\,m). The grey band marks the negligible-effect region ($|\rho| < 0.1$). While most correlations remain within the negligible band, Bandung's centrality ($\rho = +0.153$) and Semarang's POI density ($\rho = -0.193$) exceed this threshold at resolution~8, suggesting modest spatial predictor effects at the neighbourhood scale.}
\label{fig:s5}
\end{figure}

\subsection{Capacity--Congestion Scatter Plots}

\begin{figure}[!htbp]
\centering
\includegraphics[width=\textwidth]{figures/capacity_congestion_scatter.png}
\caption{Scatter plots of road capacity (lanes $\times$ speed limit) versus mean evening-peak jam factor for all three cities. Points show individual road segments; LOESS curves confirm the flat capacity--congestion relationship.}
\label{fig:s6}
\end{figure}

\subsection{Bottleneck Capacity Comparison}

\begin{figure}[!htbp]
\centering
\includegraphics[width=\textwidth]{figures/bottleneck_capacity_comparison.png}
\caption{Bottleneck capacity comparison across cities. Box plots compare the distribution of road capacity attributes (number of lanes, speed limits) between bottleneck and non-bottleneck segments identified by the graph-based capacity drop detection algorithm.}
\label{fig:s7}
\end{figure}

\subsection{Capacity Drop Spatial Distribution}

\begin{figure}[!htbp]
\centering
\includegraphics[width=\textwidth]{figures/capacity_drop_spatial_analysis.png}
\caption{Spatial distribution of capacity drop segments across all three cities. Segments are classified by whether they exhibit a statistically significant capacity drop (jam factor increase when transitioning from high-capacity to lower-capacity downstream segments).}
\label{fig:s8}
\end{figure}

\subsection{Street Orientation Analysis}

\begin{figure}[!htbp]
\centering
\includegraphics[width=\textwidth]{figures/street_orientation_polar.png}
\caption{Street orientation polar plots for Jakarta, Bandung, and Semarang. Jakarta shows a pronounced grid pattern with dominant N--S and E--W orientations; Bandung and Semarang show more heterogeneous orientations reflecting their topographic constraints.}
\label{fig:s9}
\end{figure}

%% ============================================================
\end{document}
